% =============================================================================
%  Q1 — Verification and Validation of LLM-Generated DES Experiments
%  References for this question live in refQ1.bib
%
%  GROUNDED WORKFLOW FOR THIS ANSWER:
%    1. Start from verified V&V, runtime-verification, traffic-simulation, and
%       scientific-agent sources in:
%         assets/Candidacy-Paper-reference/Question1ref/
%    2. Use extracted text only as a search aid:
%         assets/Candidacy-Paper-reference/_extracted/Question1ref/
%    3. Final claims must be grounded in:
%         (a) PDF-verified literature,
%         (b) AutoTrafficSim architecture,
%         (c) Korede's real ScalaTion/traffic-simulation experience, or
%         (d) clearly marked proposal language.
%
%  MAIN ANSWER LOGIC:
%    - Classical simulation V&V assumes a human-controlled modeling process.
%    - AutoTrafficSim changes the V&V problem because the agent proposes,
%      revises, and justifies experiments.
%    - The answer must show how invalid traces are rejected, how valid
%      surprises are preserved, what evidence makes an agent-generated
%      experiment scientifically sound, how metric overfitting is controlled,
%      and why this differs from coverage-style scenario generation.
%
%  STYLE AND DEFENSIBILITY:
%    - Use first-person technical reflection for Korede's real experience.
%    - Use proposal voice for unfinished AutoTrafficSim components.
%    - Do not claim that AutoTrafficSim has already produced results.
%    - Do not overstate what AgentGuard, AIVV, DURA-CPS, or scientific-agent
%      surveys support; keep comparisons narrow and citation-grounded.
% =============================================================================
\begin{bibunit}[plainnat]

\section{Q1.\ Verification and Validation of LLM-Generated DES Experiments}

\begin{quote}\itshape
When an LLM autonomously generates and refines discrete-event simulation experiments, part of the human researcher's traditional role as guarantor of experimental validity is transferred to the agent. Discuss the verification and validation challenges this introduces for your proposed AutoTrafficSim architecture. Specifically: Explain how your runtime invariant monitor can distinguish between a physically invalid scenario and a valid but unexpected emergent behavior worth scientific investigation. Describe what formal, empirical, and methodological evidence would be required to establish that an agent-generated experiment is scientifically sound. Discuss how the adaptive feedback loop can be prevented from overfitting to calibration metrics rather than producing genuine scientific insight. Explain what fundamentally separates your framework from large-scale automated scenario generation systems, and support your discussion with relevant literature.
\end{quote}

\subsection{Q1.1\ V\&V challenges for AutoTrafficSim}

\textit{Discuss the verification and validation challenges this introduces for your proposed AutoTrafficSim architecture.}

\citet{sargent2013introduction} and \citet{balci1998vv} describe classical
Verification and Validation (V\&V) for discrete-event simulation as a
human-organized modeling process where objectives are specified in advance, changes
are deliberate, and validity judgments are made at identifiable review points.
\citet{law1991simulation} identify manual tracing as one of the most powerful
debugging tools because it makes causality visible to the researcher. An LLM
agent may change this setting because it proposes, revises, and justifies
experiments faster than a researcher can inspect them manually.

\begin{figure}[H]
    \centering
    \includegraphics[width=\linewidth]{assets/AutoTrafficSim.png}
    \caption{Proposed AutoTrafficSim architecture. Four layers: Knowledge,
    Agent, Simulation, and Evaluation: are coupled through an Adaptive
    Feedback Loop. The Formal Validator sits inside the Simulation Layer
    (pre-simulation verification); the Runtime Invariant Monitor sits inside
    the Evaluation Layer (safety and liveness checking on the executed
    trace).}
    \label{fig:autotrafficsim}
\end{figure}

\citet{koohestani2025agentguard} identify four factors that make agentic systems
hard to verify. Each one may appears in AutoTrafficSim due to the Agentic LLM leg.

\begin{itemize}
    \item \textbf{Stochasticity and unpredictability} may be in the Agent Layer.
    \item \textbf{Hallucinations} in the Knowledge Layer and Scenario DSL.
    \item \textbf{Emergent unintended behaviors} in the Simulation Layer trace.
    \item \textbf{Susceptibility to new vulnerabilities} at pipeline boundaries.
\end{itemize}

\textbf{Stochasticity and unpredictability.}
\citet{koohestani2025agentguard} flag stochasticity as the first reason
agentic systems are hard to verify, and the same scenario prompt may produces
different runs in our own setting, so the agent becomes another source of variation sitting on top of the simulator's own randomness. When I was building my models in ScalaTion I could do manual tracing: I would open the logs, watch the animation, and walk the time-ordered \texttt{agenda} entry by entry whenever something looked off, because I was the one who set the inputs. That working style does not survive an autonomous agent that generates scenarios faster than any researcher can read them.

\textbf{Hallucinations.}
\citet{koohestani2025agentguard} place hallucinations alongside
stochasticity as a core agentic failure mode, and in AutoTrafficSim that
may appear at two specific layers: the Knowledge Layer, where the agent can
pull a parameter from the wrong context, and the Scenario DSL, where it
can quietly rewrite that parameter into something the simulator will
happily execute. AutoTrafficSim will check each proposal against the
GraphRAG retrieval source that produced it. Secondly, the agent never writes executable code on its own. It writes the Scenario DSL, and the DSL is the only thing the engine reads, similar to the same wall \citet{lopez2018sumo} put between SUMO's XML scenario format and the TraCI runtime interface so that one side cannot reach into the other.

\textbf{Emergent unintended behaviors.}
\citet{koohestani2025agentguard} treat emergent behavior as a third agentic
failure mode, and the problem in our setting is that a clean DSL check
may not protect against it. The trace may pass the gate and then produce
something that looks like a finding but is really a modeling artifact.
The proposed \textit{Runtime Invariant Monitor} will watch physical and
event-ordering invariants while the run is happening and will push
structured corrective signals back through the \textit{Adaptive Feedback
Loop}. Q1.2 carries the harder half of that question forward.

\textbf{Susceptibility to new vulnerabilities.}
\citet{koohestani2025agentguard} list new attack and failure surfaces as
the fourth agentic factor, and in AutoTrafficSim it would show up as a
pipeline problem where the output of one component becomes the input to the
next, so a small scenario-generation error could pass into the
simulator, distort the trace, skew the metric, and then be reinforced
by the feedback loop on the next iteration. \citet{rabanser2026towards}
frame the broader gap as a missing science the relaibility of AI agents.

AutoTrafficSim is not meant to remove the human researcher entirely from the
scientific process. I propose it as a way to speed up scientific knowledge
discovery in microscopic traffic simulation while leaving the harder
interpretation decisions visible. Table~\ref{tab:q11-failuremap} maps the
proposed components to the failure modes and V\&V mechanisms intended to
contain them.

\begin{table}[H]
\centering
\renewcommand{\arraystretch}{1.25}
\caption{AutoTrafficSim failure modes and proposed V\&V responses.}
\label{tab:q11-failuremap}
\begin{tabularx}{\textwidth}{
  >{\raggedright\arraybackslash}p{3.0cm}
  >{\raggedright\arraybackslash}p{3.4cm}
  >{\raggedright\arraybackslash}X}
\toprule
\textbf{Component} &
\textbf{Failure mode} &
\textbf{Proposed response} \\
\midrule
Knowledge Layer
& Wrong or unsupported retrieval
& Log the GraphRAG sources and check that they support the scenario
\citep{edge2024local}. \\
\midrule
Agent Layer
& Hypothesis drift
& Use multiple metrics and check physical invariants before accepting an
update. \\
\midrule
Scenario DSL
& Translation error
& Validate the scenario schema before it reaches the simulator. \\
\midrule
Simulation Layer
& Invalid vehicle behavior or event ordering
& Use the Runtime Invariant Monitor to reject physically invalid traces. \\
\midrule
Evaluation Layer
& Metric overfitting
& Report multiple metrics so that no single value drives refinement
\citep{manheim2019categorizing}. \\
\midrule
Adaptive Feedback Loop
& Reinforcement of a bad assumption
& Return a structured reason code when a check fails
\citep{yan2025guiding}. \\
\bottomrule
\end{tabularx}
\end{table}


% \begin{table}[H]
% \centering
% \renewcommand{\arraystretch}{1.25}
% \caption{AutoTrafficSim components, expected failure modes, and proposed V\&V checks.}
% \label{tab:q11-failuremap}
% \begin{tabularx}{\textwidth}{
%   >{\raggedright\arraybackslash}p{3.1cm}
%   >{\raggedright\arraybackslash}p{3.3cm}
%   >{\raggedright\arraybackslash}X}
% \toprule
% \textbf{Component and agentic challenge} &
% \textbf{Failure mode} &
% \textbf{Proposed V\&V response} \\
% \midrule
% Knowledge Layer\\\textit{Hallucinations}
% & Retrieval inconsistency and vocabulary mismatch
% & Log GraphRAG sources and check whether the retrieved subgraph supports the
% scenario. The Formal Validator separately checks physical plausibility
% \citep{edge2024local}. \\
% \midrule
% Agent Layer\\\textit{Stochasticity and unpredictability}
% & Hypothesis drift and objective hacking
% & Use multi-metric feedback and invariant checks before a parameter update is
% accepted. \\
% \midrule
% Scenario DSL\\\textit{Hallucinations}
% & Semantic translation errors
% & Use a schema-checked Scenario DSL to isolate agent intent from simulator
% internals \citep{balci1998vv,lopez2018sumo}. \\
% \midrule
% Simulation Layer\\\textit{Emergent unintended behaviors}
% & Scheduling-semantic mismatch
% & Check event ordering and kinematic bounds with the Runtime Invariant Monitor
% at each activation step. \\
% \midrule
% Evaluation Layer\\\textit{New vulnerabilities}
% & Metric overfitting
% & Use a multi-metric bundle so that no single metric drives refinement
% \citep{manheim2019categorizing}. \\
% \midrule
% Adaptive Feedback Loop\\\textit{New vulnerabilities}
% & Reinforcement of bad assumptions
% & Return a structured reason code when an invariant fails instead of asking the
% agent to regenerate freely \citep{yan2025guiding}. \\
% \bottomrule
% \end{tabularx}
% \end{table}


If the LLM agent proposes an invalid arrival pattern, the simulator may produce
an artificial queue. If the Evaluation Layer rewards the resulting metric, the
Adaptive Feedback Loop may reinforce the original mistake. An anomaly detector
can flag the unusual trace, but it cannot decide whether the queue is a real
traffic finding or a modeling error. This is why AutoTrafficSim needs both
runtime checks and a researcher for interpretation. This is consistent with
\citet{kwon2026aivv}, who show that flagging an anomaly is not the same as
interpreting it.

\subsection{Q1.2\ Distinguishing invalid traces from emergent behavior}

\textit{Explain how your runtime invariant monitor can distinguish between a physically invalid scenario and a valid but unexpected emergent behavior worth scientific investigation.}

For me, the invalid-versus-emergent distinction is first a physical
admissibility question, not a classification problem. If a trace violates
kinematics, it is invalid no matter how interesting it looks. If the trace
preserves the physical invariants and still produces unexpected collective
behavior, then it may be a scientific observation. Stop-and-go waves,
bottleneck breakdowns, and merge-induced oscillations can be surprising
without being invalid.

The proposed Runtime Invariant Monitor turns that distinction into a three-way
verdict:

\begin{itemize}
    \item \textbf{Reject as invalid.} The Formal Validator or the Runtime
    Invariant Monitor rejects the trace when a hard invariant fails.
    \item \textbf{Quarantine as a possible artifact.} The Runtime Invariant
    Monitor quarantines a trace whose invariants hold but whose behavior
    disappears under different random seeds, smaller time steps, or
    boundary-condition changes.
    \item \textbf{Escalate as a candidate observation.} When the invariants
    hold and the behavior is stable across those checks, the Runtime
    Invariant Monitor surfaces the trace. The researcher, not the agent, decides whether
    the result is a discovery, a known artifact, or a model defect.
\end{itemize}

The Monitor watches five invariant families:

\begin{itemize}
    \item \textbf{Flow conservation.} Vehicles entering a corridor segment
    equal vehicles exiting plus those remaining inside the segment.
    \item \textbf{Collision exclusion.}
    $x_j(t)-x_i(t)-\ell_i > 0$ for every leader $j$ and follower $i$, where
    $x$ is rear-bumper position and $\ell_i$ is the follower's length.
    \item \textbf{Kinematic bounds.}
    $0 \leq v_i(t) \leq v_{\max}$ and
    $-b_{\max} \leq a_i(t) \leq a_{\max}$
    \citep{treiber2013book}.
    \item \textbf{Lane integrity.} Vehicles do not teleport across
    non-adjacent lanes, and merge order remains consistent with queuing
    discipline.
    \item \textbf{Event ordering.} Activations are processed in timestamp order
    through ScalaTion's time-ordered \texttt{agenda}
    \citep{law1991simulation}.
\end{itemize}

\subsubsection{Why the distinction matters}

Our first lane-level validation paper showed me why both sides of this
distinction matter. We started with the Gipps model, but it could not reproduce
the acceleration profile of a steep stop-and-go regime. We moved to the
Intelligent Driver Model (IDM). The same pattern appeared with arrivals. A
Poisson process allowed two vehicles to enter the engine at the same discrete
schedule time, so we moved to a two-stage Shifted Erlang with a $\tau$
parameter enforcing a minimum inter-arrival gap. Those were invalid-scenario
catches. Later, lane-level results diverged from corridor-level results. That
mismatch was not invalid. It became the finding.

\subsubsection{Where the checks live}

The Formal Validator rejects scenarios that are already wrong before the
simulation starts. Examples include malformed DSL, densities above jam
density, arrival schedules that place two vehicles on the same discrete tick,
and lane assignments that do not match the road topology. Failed scenarios
return to the agent with a structured reason code
\citep{balci1997verification,yan2025guiding}.

The Runtime Invariant Monitor catches problems that appear only during
execution. Because simulations are stochastic, checking one run is not enough.
My proposal is to check whether the invariants continue to hold across repeated
runs and perturbations. AgentGuard provides a practical implementation
direction for threshold-based runtime checking
\citep{koohestani2025agentguard}. Zhou et al.\ provide a formal precedent for
probabilistic runtime verification under uncertainty
\citep{zhou2022runtime}. The underlying formalism is developed further in Q1.3
as part of the formal evidence pillar.

The Adaptive Feedback Loop returns failed checks to the agent for revision. It
does not decide whether an invariant-preserving anomaly is a scientific
discovery. That interpretation remains with the researcher
\citep{zanzotto2019human}. AIVV uses the same two-stage pattern: a mathematical
gate flags an anomaly, then a role-specialized council decides whether it is a
genuine fault or nuisance behavior \citep{kwon2026aivv}. Without that
researcher step, the system could suppress physically valid but surprising
traffic behavior and become a confirmation engine
\citep{treiber2013book,sugiyama2008traffic}.

\subsection{Q1.3\ Evidence for scientific soundness}

\textit{Describe what formal, empirical, and methodological evidence would be required to establish that an agent-generated experiment is scientifically sound.}

For an agent-generated experiment in AutoTrafficSim, low RMSE on a clean run
is not evidence the experiment is sound. Soundness requires three kinds of
evidence, and in the agent case all three have to be machine-checkable,
because the methodological reasoning that used to live in the researcher's
head now has nowhere to live unless the architecture writes it down.

The three pillars, mapped onto the AutoTrafficSim components in
Figure~\ref{fig:autotrafficsim}, are:

\begin{itemize}
    \item \textbf{Formal evidence.} The Formal Validator in the Simulation
    Layer confirms each scenario satisfies its DSL contract before the run,
    and the Runtime Invariant Monitor in the Evaluation Layer confirms the
    trace preserves the five physical invariants from Q1.2 during the run.
    \item \textbf{Empirical evidence.} Metrics computation in the Evaluation
    Layer compares simulated flow, speed, occupancy, and travel time against
    external ground truth such as PeMS, and reports whether the conclusion
    still holds across the scenario sensitivity sweep driven by the Agent
    Layer.
    \item \textbf{Methodological evidence.} The Knowledge Layer records the
    GraphRAG retrieval trace behind each assumption, and the Agent Layer
    records the hypothesis and revision lineage that produced each scenario.
\end{itemize}

The three-pillar framing is not new. It is my translation of standard
simulation V\&V into the agent setting. \citet{sargent2010verification}
separates conceptual validity, operational validity, and computerized-model
verification. Computerized-model verification becomes the formal pillar.
Operational validity becomes the empirical pillar. Conceptual validity
becomes the methodological pillar, because the agent has to record why each
modeling choice was made. \citet{law1991simulation} adds a constraint I need
to respect across all three: simulation validation is usually about closeness
for a stated purpose, not proving that simulated and real outputs are
identical.

The pillars cut across the six layers of Table~\ref{tab:q11-failuremap}
rather than lining up one-to-one with them. The Adaptive Feedback Loop
threads all three together by carrying diagnostic signal back into the next
iteration.

Verifying code, objectives, and runtime behavior in ScalaTion by hand made
this concrete for me. I was asking formal questions, such as whether an event
fired in causal order. I was asking empirical questions, such as whether the
flow profile matched what IDM should produce at that demand level. I was also
asking methodological questions, such as why we chose the Shifted Erlang
$\tau$ for the arrival process. The formal and empirical questions left
traces in logs and metrics. The methodological questions often lived in my
head. Once the agent generates scenarios at machine speed, that reasoning
disappears unless AutoTrafficSim records it. Surveys of scientific agents
make the same point in broader terms: reproducibility and procedural
traceability cannot be treated as afterthoughts
\citep{zhou2025autonomous,ren2503towards}.

\subsubsection{Where each pillar is produced}

\textbf{Formal evidence.} The Formal Validator and the Runtime Invariant
Monitor would produce records that can be inspected after each experiment.
These include a DSL contract-satisfaction log from the Formal Validator,
invariant-preservation and flow-conservation records from the Runtime
Invariant Monitor, and timestamp-ordered event-activation records from the
ScalaTion simulation engine's agenda. The monitor is still a proposal. My
first step is simple. The simulator records whether the physical invariants
hold across repeated runs. One clean run is not enough because random seeds,
the Shifted Erlang arrival process, and probabilistic lane-change decisions
can produce different traces from the same scenario. AgentGuard gives me a
practical direction for the next step. It observes agent traces, learns an
Agentic Markov Decision Process (AMDP) online, and uses Probabilistic Model
Checking (PMC) to trigger alerts or automated responses when a threshold is
breached \citep{koohestani2025agentguard}. I am not claiming that
AutoTrafficSim already has this backend. My proposal is to adapt the
AgentGuard pattern so the workflow can also monitor the agent's revisions
and identify when intervention is needed.

\textbf{Empirical evidence.} The empirical checks depend on the experiment.
For ordinary freeway validation, AutoTrafficSim would compare simulated
flow, speed, occupancy, and travel time against observed traffic data such
as PeMS. My earlier lane-level work showed why aggregate fit is not enough.
A corridor-level match can still hide a lane-level problem. For the wildfire
evacuation case in Q2, the evaluation has to go beyond ordinary traffic fit.
The framework would report evacuation clearance time, exposed vehicles,
queue duration, bottleneck location, and variation across repeated runs.
Those outcomes would be compared across the scenario matrix of fire spread,
departure timing, route choice, contraflow, phased evacuation, and household
vehicle demand. I am not trying to find one best evacuation trace. I want to
know which assumptions change the decision, and whether the conclusion still
holds when those assumptions change. Parameter sensitivity analysis
completes the pillar \citep{sargent2010verification,law1991simulation}. A
result is more useful when the conclusion holds across reasonable
assumptions, and more honest when the framework shows where it does not.

\textbf{Methodological evidence.} AutoTrafficSim would record where each
experiment came from and why it changed. The GraphRAG retrieval trace would
show which knowledge-graph sources the agent consulted
\citep{edge2024local}. The hypothesis lineage would show how a candidate
came from earlier results. Parameter-change records would connect each
revision to the invariant violation, metric shift, or anomaly that triggered
it. Reproducibility logs would preserve seeds, DSL versions, validator
commit hashes, and metric outputs so another researcher can reconstruct the
run. Without these records, an autonomous experiment cannot be distinguished
from a search process that happened to find a good metric value
\citep{rabanser2026towards}. Without them, the researcher who has to
interpret an escalated trace from Q1.2 has nothing to work from.

\subsection{Q1.4\ Anti-overfitting in the Adaptive Feedback Loop}

\textit{Discuss how the adaptive feedback loop can be prevented from overfitting to calibration metrics rather than producing genuine scientific insight.}

To keep the Adaptive Feedback Loop from overfitting to calibration metrics,
AutoTrafficSim relies on three components already shown in
Figure~\ref{fig:autotrafficsim}:

\begin{itemize}
    \item Metrics Computation will report more than one metric.
    \item The Runtime Invariant Monitor will reject parameter sets that
    violate physical constraints.
    \item The Adaptive Feedback Loop will return the failed check to the
    agent rather than reward an RMSE improvement alone.
\end{itemize}

This keeps a low aggregate error from hiding an invalid microscopic run. I saw the underlying problem during manual calibration of the ScalaTion
microscopic traffic simulator, before the AutoTrafficSim Agent Layer
existed. Using Simultaneous Perturbation Stochastic Approximation (SPSA) on
IDM parameters, the optimizer returned parameter sets with a good RMSE
while the microscopic vehicle trajectories were still physically
inadmissible. Vehicles were spaced closer than physically possible or
accelerated beyond kinematic limits. Optimizing the metric alone is
therefore not enough. \citet{manheim2019categorizing} formalize this as
Goodhart's Law. A GraphRAG-driven Agent Layer proposing parameter
revisions at machine speed would reproduce that failure faster than I did
by hand.

A later extension could add clone-and-promote gating from AIVV, where an
update is tested on a temporary clone before promotion
\citep{kwon2026aivv}. I would treat that as an extension, not as a
completed part of the current architecture.

\subsection{Q1.5\ Differentiation from automated scenario generation}

\textit{Explain what fundamentally separates your framework from large-scale automated scenario generation systems, and support your discussion with relevant literature.}

I can already run the ScalaTion engine on hundreds of parameter
combinations. That gives coverage. It does not tell me whether any run
answered a useful scientific question, whether a result is physically
admissible, whether the conclusion depends on one fragile assumption, or
why the next experiment should be run.

What separates AutoTrafficSim from a large-scale automated scenario
generator is therefore not throughput but a checked and traceable path
from an assumption, to an experiment, to a result a researcher can
defend. Before the run, the Knowledge Layer records where the assumption
came from, the Agent Layer turns it into an experiment, and the Formal
Validator rejects a scenario that is already physically inadmissible.
After the run, the Runtime Invariant Monitor checks the trace, Metrics
Computation reports what changed, and the Adaptive Feedback Loop returns
that evidence to the agent before another experiment is proposed. If the
trace is physically valid but still surprising, the result is sent to
the researcher for interpretation rather than silently discarded.

Existing tools such as SUMO, Scenic, CARLA, and OpenSCENARIO are useful
for defining and executing scenarios
\citep{lopez2018sumo,fremont2019scenic,dosovitskiy2017carla,asam2024openscenario}.
I am not proposing AutoTrafficSim as a replacement for those tools or
for ScalaTion. The simulation engine remains the engine. AutoTrafficSim
is the agentic layer around it.

This matters for the wildfire evacuation case in Q2. A large scenario
generator can vary fire spread, departure timing, rerouting, phased
evacuation, household vehicle demand, and contraflow. AutoTrafficSim
goes beyond enumerating those combinations: it shows which assumption
caused the conclusion to change and whether the result still holds when
reasonable assumptions change.

Recent surveys treat agentic scientific experimentation as an emerging
direction because it connects hypothesis generation, experiment
execution, and result refinement in one system
\citep{ren2503towards,zhao2025llm,zhou2025autonomous}.
\citet{zhou2025autonomous} describe autonomous scientific discovery as a
loop of hypothesis discovery, experimental design and execution, and
result analysis and refinement. That is the idea I am trying to make
concrete for traffic simulation.


\bigskip
\phantomsection
\addcontentsline{toc}{subsection}{References}
\putbib[refQ1]
\end{bibunit}
